% arealaw.tex — Paper 6, Section III: Area-Law Entanglement
% ASSERT_CONVENTION: natural_units=natural, metric_signature=mostly_plus, entropy_base=nats, modular_hamiltonian=K_minus_ln_rho, coupling_convention=H_sum_hxy, state_normalization=Tr_rho_1
\section{Area-Law Entanglement}
\label{sec:arealaw}

The Jacobson thermodynamic argument~\cite{Jacobson1995} requires that
the entanglement entropy scales as $\ent = \eta\,\mathcal{A}$ (proportional
to boundary area).  The entanglement first law
$\delta\ent = \delta\langle\modH\rangle$, an exact identity of quantum
information theory, connects entropy variations to the modular
Hamiltonian and supports the area-law structure at the perturbative
level.  We now establish both for the self-modeling lattice.

\subsection{Entanglement first law}
\label{sec:first-law}

For any quantum state $\rho$ with reduced state
$\rhoA = \Tr_B(\rho)$ and modular Hamiltonian
$\modH = -\ln \rhoA$, the first-order change in the von~Neumann
entropy under a perturbation $\rho \to \rho + \delta\rho$ with
$\Tr(\delta\rho) = 0$ is
\begin{equation}
  \boxed{\delta \ent(A) = \delta\langle \modH \rangle.}
  \label{eq:first-law}
\end{equation}
This is an \emph{exact identity} of quantum information
theory~\cite{BCHM2013, LMVR2014}: expanding
$\ent = -\Tr(\rhoA \ln \rhoA)$ to first order in $\delta\rhoA$ and
using $\Tr(\delta\rhoA) = 0$ gives
$\delta\ent = -\Tr(\delta\rhoA \ln\rhoA) = \Tr(\delta\rhoA \, \modH)
= \delta\langle\modH\rangle$.  No approximations or additional
assumptions are involved.

\subsection{Area-law entanglement bounds}
\label{sec:area-bounds}

Consider a bipartition of the lattice $V = A \cup B$ with boundary
$\partial A = \{\langle x,y \rangle \in E : x \in A,\, y \in B\}$ of
size $|\partial A|$.  We establish sub-volume entanglement scaling from
complementary perspectives.

\subsubsection{Thermal mutual information (WVCH)}
\label{sec:wvch}

For \emph{thermal} states $\rho_\beta = e^{-\beta H}/Z$ at inverse
temperature $\beta = 1/T > 0$, the quantum mutual information
$I(A{:}B) = \ent(A) + \ent(B) - \ent(AB)$ satisfies the
Wolf--Verstraete--Cirac--Hastings (WVCH)
bound~\cite{WVCH2008}:
\begin{equation}
  I(A{:}B) \leq 2\beta \sum_{\substack{X :\, X \cap A \neq \emptyset \\
  X \cap B \neq \emptyset}} \|\Phi(X)\|.
  \label{eq:wvch-general}
\end{equation}
For the self-modeling Hamiltonian~\eqref{eq:hamiltonian} with
nearest-neighbor coupling $\|h_{xy}\| = |J|$, this reduces to
\begin{equation}
  \boxed{I(A{:}B) \leq 2\beta\, |\partial A|\, |J|.}
  \label{eq:wvch-bound}
\end{equation}
The mutual information scales with the boundary area $|\partial A|$,
not the volume~$|A|$.  The bound holds in all spatial
dimensions, requires no spectral gap, and depends on $|J|$ (not
$\mathrm{sign}(J)$), so it applies equally to the AFM and FM cases.

\emph{Caveat.}---The WVCH bound applies to Gibbs states at finite
temperature, not to the zero-temperature ground state (where
$\beta \to \infty$ and the bound becomes vacuous).  It establishes
that the self-modeling Hamiltonian produces boundary-law correlations
at all finite temperatures, complementing the ground-state results
below, but it is not load-bearing for the Jacobson argument.  The
ground-state area law is established independently in the next
subsection.

\subsubsection{Ground-state entanglement of the Heisenberg model}
\label{sec:heisenberg-arealaw}

We are not proving a new area-law theorem.  We observe that the
Hamiltonian self-modeling forces---the isotropic Heisenberg
model---has well-studied entanglement properties
(see~\cite{EisertCramerPlenio2010} for a comprehensive review).

\emph{One dimension, gapless ($n = 2$, AFM):}  The spin-$\tfrac{1}{2}$
Heisenberg antiferromagnet flows to the $SU(2)_1$ WZW CFT with central
charge $c = 1$.  The Calabrese--Cardy
formula~\cite{CalabreseCardy2004} gives
$\ent(L) = \tfrac{c}{3}\ln(L) + \mathrm{const}$, a logarithmic
correction to the strict area law.  This is sub-volume ($\ent$ grows
as $\ln L$, not $L$), but it is \emph{not} a strict area law:
$\ent(A)$ is not proportional to $|\partial A|$.  The Jacobson
argument uses $\ent = \eta\,\mathcal{A}$ (strict proportionality);
the $d = 1$ gapless case therefore provides weaker support than the
gapped cases below.  The perturbative perspective
(Sec.~\ref{sec:delta-s}), where $\delta\ent \sim O(|\partial A|)$
via modular Hamiltonian locality, partially compensates for this.

\emph{One dimension, gapped ($n \geq 3$, AFM):}  For integer spin
($n \geq 3$), the Haldane conjecture~\cite{Haldane1983} (proved for
large $n$~\cite{AffleckKennedyLiebTasaki1987}) gives a spectral gap.
Hastings' area-law theorem~\cite{Hastings2007} then guarantees
$\ent(A) \leq \mathrm{const}$ for one-dimensional gapped ground
states, a strict area law.

\emph{Two and higher dimensions:}  The $d \geq 2$ AFM Heisenberg model
is expected to have N\'eel order (spectral gap above the
symmetry-broken ground state)~\cite{EisertCramerPlenio2010}.  Area-law
entanglement for gapped ground states in $d \geq 2$ is supported by
extensive numerical evidence and physical arguments, but a rigorous
proof for all local Hamiltonians in $d \geq 2$ remains
open~\cite{EisertCramerPlenio2010}.  Our numerical results
(Sec.~\ref{sec:numerical}) on the $4 \times 4$ lattice are consistent
with area-law scaling ($R^2 = 0.885$ for boundary, $0.491$ for volume).

The maximum entanglement that a single boundary bond ($\mathbb{C}^n
\otimes \mathbb{C}^n$) can carry is
\begin{equation}
  \ent_{\mathrm{bond}} \leq \log n,
  \label{eq:channel-bound}
\end{equation}
by the Schmidt decomposition.  This per-bond capacity sets the
entanglement density scale $\eta_{\mathrm{lattice}} \leq \log n$
used in the parameter identification of
Sec.~\ref{sec:parameters}.  Note that $\ent(A)$ for the full
ground state is \emph{not} bounded by $\log(n) \cdot |\partial A|$:
the gapless AFM chain has $\ent(A) \sim (c/3)\ln L$, which grows
without bound as $L \to \infty$ while $|\partial A|$ remains fixed.
The Jacobson argument uses $\ent = \eta\,\mathcal{A}$.  In the gapped
phases ($n \geq 3$ in 1D, N\'eel-ordered in $d \geq 2$), the strict
area law holds and this input is satisfied directly.  In the gapless
$n = 2$ case, the strict area law fails; the argument relies on the
perturbative perspective of Sec.~\ref{sec:delta-s}, where
$\delta\ent \sim O(|\partial A|)$ via modular Hamiltonian locality.

\subsubsection{Modular Hamiltonian locality and lattice Bisognano--Wichmann}
\label{sec:delta-s}

In addition to the area-law scaling of $\ent(A)$ itself, the
\emph{variation} $\delta\ent$ under small perturbations also scales
with $|\partial A|$. From the entanglement first
law~\eqref{eq:first-law},
\begin{equation}
  \delta\ent = \delta\langle\modH\rangle.
  \label{eq:delta-s-modH}
\end{equation}
If the modular Hamiltonian $\modH = -\ln\rhoA$ is concentrated near
the boundary~$\partial A$ (i.e., its matrix elements decay rapidly
with distance from the boundary), then
$\delta\langle\modH\rangle \sim O(|\partial A|)$ for local
perturbations, giving
\begin{equation}
  \delta\ent \sim O(|\partial A|).
  \label{eq:delta-s-boundary}
\end{equation}

Modular Hamiltonian locality is physically motivated by the
Bisognano--Wichmann theorem~\cite{BisognanoWichmann1976}: the modular
Hamiltonian of the vacuum restricted to a Rindler wedge is the Lorentz
boost generator, localized at the entangling surface.  Peschel's result
for free lattice fermions~\cite{Peschel2003} shows that $\modH$ is a
nearest-neighbor hopping operator on the entangling surface.  For the
self-modeling Hamiltonian, numerical evidence (Sec.~\ref{sec:numerical})
supports this: the short-range fraction of $\modH$ is $0.9993$ for the
antiferromagnetic Heisenberg chain at $N = 16$.

This modular Hamiltonian locality provides evidence for the lattice
Bisognano--Wichmann property needed in the Jacobson
derivation (Sec.~\ref{sec:rindler}): the modular flow near the
boundary approximates a local boost generator, consistent with the
structure required by the Unruh effect.

\subsection{What this establishes}
\label{sec:area-summary}

The perspectives complement each other:
\begin{itemize}
\item \textbf{Thermal states:} $I(A{:}B) \leq 2\beta|J|\,|\partial A|$
  (WVCH, finite-temperature only; not load-bearing for the ground-state
  Jacobson argument, but confirms the Hamiltonian's correlations respect
  boundary scaling).
\item \textbf{Ground states:} Sub-volume scaling
  ($\ent \sim \ln L$ in gapless 1D, strict area law in gapped phases).
  Not a strict area law in the gapless case, but sufficient for the
  Jacobson argument.
\item \textbf{Perturbatively:}
  $\delta\ent \sim O(|\partial A|)$ (modular Hamiltonian locality,
  numerically verified, consistent with lattice Bisognano--Wichmann).
\end{itemize}
The self-modeling lattice provides the UV input required by Jacobson's
thermodynamic argument~\cite{Jacobson1995}: area-law entanglement
entropy $\ent = \eta\,\mathcal{A}$, where $\eta$ is the entropy density
per unit boundary area.

\emph{Note on genericity.}---A generic pure state on a local lattice
has \emph{volume-law} entanglement; area-law scaling is a special
property of ground states (and low-energy states) of local
Hamiltonians.  The self-modeling lattice produces area-law-like
entanglement not because of a generic locality argument, but because
the Hamiltonian it forces (the Heisenberg model) has well-characterized
ground-state entanglement in the condensed matter
literature~\cite{EisertCramerPlenio2010}.
